\section{Software Product Vision}

FREVIA is a web-based freelance marketplace built around a single proposition: an engagement between two strangers should be governable from beginning to end without either party having to trust the other on faith. The product therefore does not compete on the size of its talent pool. It competes on the completeness of the record it produces for each engagement -- a verified identity, a signed contract, a funded milestone, a reviewed deliverable, and, when the parties disagree, a documented arbitration decision.

The primary users are clients who need a bounded, verifiable engagement and freelancers who need to know that payment exists before work begins. A third role, the expert, provides a public professional profile with an experience and certification record, intended for users whose value lies in review and consultation rather than in delivery. A fourth role, the administrator, is an operator rather than a customer: it approves profile revisions, reviews identity documents, orders the moderation queue, resolves disputes, and adjusts what each role is permitted to do. The same person may hold more than one role, and the interface follows the role that is currently active in the session.

For the client, the vision is a hiring flow that never requires a leap of faith. Posting a job captures a fixed budget range, a deadline, the required skills, the categories, and the number of positions to fill. Proposals arrive with the freelancer's profile, skill proficiency, and attachments, and can be compared side by side before a candidate is hired. Hiring generates a contract that both parties must sign; funding a milestone moves the amount into escrow; approving a submission completes the milestone and makes the funded amount releasable, after which the client releases it as a payout and settles the fixed platform fee. At any point the client can see the money state of the contract, the status of each milestone, and the files that have been exchanged.

For the freelancer, the vision is the mirror image. A profile is built from structured evidence: a title, education, certifications, languages, work experience, skills with a proficiency level, portfolio items with technologies and a project link, and a curriculum vitae that can be parsed to propose skills automatically. Identity verification turns a claim into a badge. Job discovery is driven by explicit filters and saved searches rather than by an opaque ranking, bookmarks keep interesting postings within reach, and proposals can be prepared as drafts, submitted with attachments and an expiry date, edited, or withdrawn. Once hired, the freelancer can see which milestones are funded, submit work with files and links, respond to a change request with a due date, and escalate to a dispute if the review does not converge.

The organising principle behind these two journeys is that every state transition must be visible to both sides and must be reversible only through a recorded action. A milestone is not silently completed; it is submitted, optionally returned for changes, and then approved. A contract is not implicitly accepted; both signatures are stored. A dispute is not resolved by an email; the split, the reasoning, and the reviewing administrator are all persisted. This is why the platform invests in a transaction ledger, a notification record, and an explicit status model rather than in a chat-first experience.

Artificial intelligence supports this principle rather than replacing it. The AI components in FREVIA are narrow by design: a CV analyser that proposes skills with a proficiency estimate, an identity verification pipeline that reads a document, matches a face, and checks liveness, and a forum moderation step that scores content before a human approves it. None of them makes a decision that binds a user. The client still chooses whom to hire, the administrator still approves or rejects, and the freelancer still confirms what goes onto the profile. Keeping a human at each decision point is what allows the platform to explain any outcome to the person affected by it.

The technical vision follows the same logic. FREVIA is a monorepo that holds a NestJS API, a Next.js web client, a shared package of validation schemas and generated types, and a shared user-interface package. Because the validation schemas are shared, a rule such as a minimum proposal length or a maximum milestone amount is enforced identically in the browser and on the server. Access control is data-driven: permissions are rows that bind an HTTP method and a path to a module, roles are granted permissions, and a single guard checks every request against the role active in the session, which means authorisation can be corrected in production data rather than in code. Money is handled exclusively through Stripe -- Express onboarding for payouts, PaymentIntents and hosted Checkout Sessions for funding, transfers for release, refunds for reversal, and a signature-verified webhook for asynchronous state. The system is containerised so that a fresh environment can be started with one command.

The longer-term direction is to deepen rather than to widen. The most valuable extensions are the ones that reuse the existing record: additional verification methods behind the same document and liveness interface, hourly contracts alongside the current fixed-price model, multi-currency settlement once Stripe support allows it, and richer reputation signals derived from completed contracts rather than from self-reported claims. Each of these extends a structure that already exists, which is the practical test the team applies when deciding what belongs on the roadmap.
